\documentclass[preprint,review,12pt,authoryear]{elsarticle}

% Core math and symbols
\usepackage{amssymb}
\usepackage{amsmath}

% Tables and formatting
\usepackage[T1]{fontenc}
\usepackage{multirow}
\usepackage{array}
\usepackage{tabularx}
\usepackage{booktabs}
\usepackage{makecell}
\usepackage{xurl}
\usepackage{placeins}
\usepackage{lineno}
\renewcommand\linenumberfont{\normalfont\tiny\sffamily}
\setlength\linenumbersep{3mm}
\modulolinenumbers[1]

% Graphics path (figures are in the project root)
\usepackage{graphicx}
\graphicspath{{../}}

% Float placement parameters (kept from source)
\renewcommand{\topfraction}{0.9}
\renewcommand{\bottomfraction}{0.8}
\renewcommand{\textfraction}{0.1}
\renewcommand{\floatpagefraction}{0.8}
\setcounter{topnumber}{4}
\setcounter{bottomnumber}{3}
\setcounter{totalnumber}{6}

% Hyperlinks for cross-references
\usepackage[hidelinks]{hyperref}

% Journal (optional placeholder)
\journal{Remote Sensing of Environment}

\begin{document}

\begin{frontmatter}

% Title
\title{Resolving signal ambiguity in earthquake precursor detection with an environment-specific deep learning framework for AMSR-2 data}

% Authors and affiliations (mapped from original superscripts)
\author[inst1]{Pan Xiong\corref{cor1}}
\author[inst2]{Cheng Long}
\author[inst3]{Huiyu Zhou}
\author[inst4,inst5]{Roberto Battiston}
\author[inst6]{Angelo De Santis}
\author[inst7]{Xuhui Shen\corref{cor1}}

\affiliation[inst1]{organization={Institute of Earthquake Forecasting, China Earthquake Administration},
            city={Beijing},
            country={China}}
\affiliation[inst2]{organization={College of Computing and Data Science, Nanyang Technological University},
            city={Singapore},
            country={Singapore}}
\affiliation[inst3]{organization={School of Computing and Mathematical Sciences, University of Leicester},
            city={Leicester},
            country={United Kingdom}}
\affiliation[inst4]{organization={Department of Physics, University of Trento},
            city={Trento},
            country={Italy}}
\affiliation[inst5]{organization={National Institute for Nuclear Physics, The Trento Institute for Fundamental Physics and Applications},
            city={Trento},
            country={Italy}}
\affiliation[inst6]{organization={Istituto Nazionale di Geofisica e Vulcanologia},
            city={Rome},
            country={Italy}}
\affiliation[inst7]{organization={National Space Science Center, Chinese Academy of Sciences},
            city={Beijing},
            country={China}}

\cortext[cor1]{Corresponding authors.}

% Abstract
\begin{abstract}
Satellite-based earthquake monitoring using microwave brightness temperature (MBT) is fundamentally challenged by the signal ambiguity problem, where pre-seismic anomalies are difficult to distinguish from non-seismic environmental variations. We address this challenge by developing an environment-specific, knowledge-guided deep learning framework to analyze Advanced Microwave Scanning Radiometer 2 (AMSR-2) data. Our approach first classifies the Earth's surface into five distinct environmental zones. Within each zone, we use association rule mining to discover reliable precursor signatures (e.g., specific 89 GHz H-polarization patterns over marine environments). This knowledge is then integrated into a novel Weight-Enhanced Feature-Tailored Transformer (WE-FTT) model to guide its focus. Validated on a dataset of 346.56 million MBT measurements across 154 major earthquakes (M$\ge$7.0, 2013-2023), the WE-FTT achieves a Matthews Correlation Coefficient of approximately 0.84, significantly outperforming conventional global methods. This environment-specific approach provides a robust methodological foundation for developing future operational earthquake early warning systems.
\end{abstract}

% Keywords
\begin{keyword}
AMSR-2 \sep Microwave brightness temperature \sep Earthquake precursor detection \sep Environment-specific analysis \sep Knowledge-guided deep learning \sep Weight-Enhanced Feature-Tailored Transformer
\end{keyword}

\end{frontmatter}

% Enable line numbers for main text
% \linenumbers

\section{Introduction} 
Satellite microwave remote sensing, particularly using passive radiometry in the 6.9--89 GHz range, offers unique capabilities for monitoring pre-earthquake surface anomalies due to its all-weather operation and sensitivity to dielectric property variations \citep{pulinetsLithosphereAtmosphereIonosphere2011,wuGEOSSbasedThermalParameters2012}. The Advanced Microwave Scanning Radiometer 2 (AMSR-2) provides global coverage with multi-frequency brightness temperature (MBT) observations that potentially capture stress-induced electromagnetic emissions from the lithosphere-atmosphere coupling process \citep{kawanishiAdvancedMicrowaveScanning2003}. However, the operational implementation of this technology faces a fundamental challenge: distinguishing genuine seismic precursors from environmental variations in MBT data---a problem commonly referred to as signal ambiguity \citep{troninRemoteSensingEarthquakes2006,tramutoliRobustSatelliteTechniques2013}.

The signal ambiguity problem manifests as similar MBT anomaly patterns arising from both seismic and non-seismic sources. Environmental factors such as soil moisture fluctuations, vegetation phenology, and atmospheric water vapor can produce MBT signatures that mimic those potentially associated with earthquake preparation \citep{qiCharacteristicBackgroundMicrowave2023}. While previous studies have documented promising anomalies preceding individual earthquakes \citep{akhoondzadehMultiPrecursorsAnalysis2018,qiSpatiotemporallyWeightedTwostep2020}, systematic global analyses have failed to establish statistically robust detection thresholds with acceptable false positive rates \citep{gellerEarthquakesCannotPredicted1997,cuiSatelliteThermalInfrared2019,contiSystematicReviewMetaAnalysis2021,ciceroneSystematicCompilationEarthquake2009}. This inconsistency indicates that prevailing methodologies, which often treat Earth's surface as homogeneous, may be fundamentally misaligned with the physical nature of pre-seismic signals.

We hypothesize that this limitation stems from a critical oversight: the failure to account for the diverse microwave radiative transfer characteristics across different environments \citep{njokuRetrievalLandSurface1999,guptaMicrowaveEmissionScattering2014}. A universal detection algorithm inevitably conflates environment-specific signatures with potential seismic signals, amplifying rather than resolving ambiguity. 

To address this, we develop a novel framework that integrates environmental stratification with knowledge-guided deep learning, leveraging recent advances in machine learning for Earth observation \citep{reichsteinDeepLearningEarth2019,zhuDeepLearningRemote2017}. This study makes three primary contributions to remote sensing methodology. First, we implement an automated surface classification that partitions Earth into five distinct zones based on vegetation and soil moisture, enabling an analysis that respects the underlying physics of microwave emission. Second, we introduce the Weight-Enhanced Feature-Tailored Transformer (WE-FTT), a new deep learning architecture that explicitly incorporates domain knowledge of environment-specific precursor patterns derived from association rule mining. Third, we provide comprehensive validation using 346.56 million AMSR-2 MBT measurements for 154 major (M$\ge$7.0) earthquakes (2013--2023), demonstrating that our approach achieves a Matthews Correlation Coefficient of approximately 0.84, a substantial improvement over conventional global methods.

The overarching goal of this research is to quantitatively demonstrate that reliable MBT-based earthquake precursors can be identified when their environment-specific characteristics are properly modeled. This work provides a robust methodological foundation for developing more reliable satellite-based seismic hazard assessment systems by resolving the long-standing signal ambiguity problem.

The remainder of this paper is organized as follows: Section 2 describes the data sources and the methodological framework, including the environment-specific analysis and the WE-FTT model. Section 3 presents the empirical results. Section 4 discusses the findings and their implications, and Section 5 provides the conclusions.

\section{Data and Data Preprocessing}
\section{Data and Methods}\label{sec:data_methods}

The methodological framework (Figure~\ref{fig:fig1}) integrates data acquisition, environment-specific partitioning, knowledge mining, and deep learning to address the signal ambiguity problem in earthquake precursor detection. Each component of this framework is detailed in the following subsections.

\subsection{Data Acquisition and Preprocessing}\label{sec:data_acquisition}

The primary dataset comprises microwave brightness temperature (MBT) measurements from the Advanced Microwave Scanning Radiometer 2 (AMSR-2) instrument, operational since 2012 \citep{kasaharaStatusAMSR22012}. AMSR-2 provides observations across six frequency bands (6.9 to 89 GHz) in both vertical and horizontal polarizations. This study utilizes ten of these channels (6.9, 10.65, 23.8, 36.5, and 89.0 GHz for both H and V polarizations). The instrument's spatial resolution varies with frequency, from approximately 3 $\times$ 5 km to 24 $\times$ 42 km \citep{kachiStatusGCOMW1AMSR22008}. To minimize interference from diurnal solar radiation and anthropogenic activity, we use only nighttime (descending-mode) data, which provides a more stable baseline for anomaly detection \citep{maedaDiscriminationLocalFaint2008}. The entire dataset, spanning January 2013 to August 2023, was gridded to a 0.25$^{\circ}$ resolution for consistency with auxiliary data and global analysis practices \citep{njokuRetrievalLandSurface1999,wuIdentifyingSeismicAnomalies2024}. The multi-frequency nature of AMSR-2 is particularly advantageous, as lower frequencies (e.g., 6.9 GHz) provide deeper penetration into the soil column, while higher frequencies (e.g., 89.0 GHz) are more sensitive to surface skin temperature and atmospheric effects, enabling a quasi-tomographic analysis of potential pre-seismic phenomena.

To effectively isolate potential seismic signals, we integrated several auxiliary datasets. Daily soil moisture data were sourced from the NASA GES DISC (0.25$^{\circ}$ resolution) \citep{fangGlobalLandData2009}, and daily vegetation coverage data were derived from the ERA5 reanalysis dataset (0.25$^{\circ}$ resolution) \citep{hersbachERA5GlobalReanalysis2020,munoz-sabaterERA5landStateoftheartGlobal2021}. For ground truth, we used the global earthquake catalog from the U.S. Geological Survey (USGS), which includes 32,123 events with magnitudes M$\ge$4.8 from January 1, 2013, to August 1, 2023. Although our initial analysis considered this full range, consistent MBT anomalies were found to be predominantly associated with major earthquakes (M$\ge$7.0). Therefore, this study focuses its primary analysis on these larger events to develop a robust methodological framework.

Quality control procedures were applied to the AMSR-2 data, including the removal of pixels affected by Radio Frequency Interference (RFI), cloud screening based on the 89 GHz polarization difference, and interpolation of missing values using bilinear methods where data gaps were less than three consecutive pixels. AMSR-2 fill values (e.g., 655.35 K) were excluded from all statistical analyses and model training.

\subsection{Environment-Specific Analysis Framework}\label{sec:env_framework}

\subsubsection{Surface Type Classification}\label{sec:surface_classification}
Given the high sensitivity of MBT signals to surface conditions, a crucial first step is to analyze data within homogeneous environmental contexts. We classified the global surface into five distinct zones based on vegetation coverage, soil moisture, and land-sea boundaries, with the specific criteria summarized in Table~\ref{tab:table2}. These five zones were designed to represent fundamentally different microwave radiative transfer regimes, which are primarily governed by surface water content and vegetation structure. For implementation, daily vegetation and soil moisture data were processed to our 0.25$^{\circ}$ analysis grid. After applying a land-sea mask to identify Zone A (Marine), terrestrial regions were classified into Zones B--E based on the defined thresholds. This classification enables the development of tailored analytical approaches for each zone's unique characteristics \citep{oweMethodologySurfaceSoil2001}.

\subsubsection{Seismological Controls}\label{sec:seismo_controls}
To address the valid point on seismological factors, we introduced a second layer of stratification based on source parameters. In addition to the environmental zoning, we stratified all earthquake-related samples by magnitude bins (7.0--7.4, 7.5--7.9, and $\ge$8.0) and focal depth bins (0--70 km, 70--150 km, and >150 km). This allows us to control for source-side variability and enables a more direct comparison with other precursor studies, such as those based on TEC, that have reported dependencies on these parameters \citep{Shah2020TotalEC}. The results of this stratified analysis are reported as conditional MCC and FPR values in the supplementary material (Table~\ref{tab:tableS3}).

\subsubsection{Identification of Seismic and Non-seismic Data}\label{sec:data_identification}
To construct our training and evaluation datasets, we first identified MBT measurements associated with seismic events. For each of the 154 major earthquakes (M$\ge$7.0) in our catalog, we defined a spatiotemporal window of influence. The global distribution of these earthquakes across the five defined environmental zones is visualized in Figure~\ref{fig:fig11}.The temporal window spanned from 20 days before to 10 days after each event to capture both potential short-term precursors and post-seismic effects \citep{akhoondzadehMultiPrecursorsAnalysis2018,qiSpatiotemporallyWeightedTwostep2020}. The spatial extent was defined by the Dobrovolsky radius (R = $10^{0.43M}$ km), a standard empirical relationship for the earthquake preparation zone \citep{dobrovolskyEstimationSizeEarthquake1979}. For marine events (Zone A), only shallow earthquakes (focal depth < 70 km) were included, and we applied additional controls for potential artifacts by excluding data from within ±48 hours of a documented tsunami arrival and filtering out days with extreme sea states. Any MBT measurement falling within the spatiotemporal window of a qualifying earthquake was labeled as "earthquake-related."

To create a balanced dataset for supervised learning, we then generated a corresponding non-seismic dataset for each of the five environmental zones. This was achieved by randomly sampling a number of non-event MBT measurements equal to the count of the earthquake-related samples for that zone. This procedure ensures that each zone-specific dataset is balanced, minimizing potential classification bias while preserving unique environmental characteristics.

\subsection{Knowledge Mining of Precursor Signatures}\label{sec:knowledge_mining}
To move beyond "black-box" detection, we implemented a data mining approach to identify consistent MBT patterns that reliably precede earthquakes within each environmental context.

\subsubsection{Clustering Analysis}\label{sec:clustering}
As a prerequisite for association rule mining, we first discretized the continuous MBT data. We chose the K-means algorithm for its computational efficiency with large-scale datasets \citep{hanConceptsTechniques2006}. The optimal number of clusters (K) for each of the 10 AMSR-2 channels was determined using the Elbow Method, which analyzes the Sum of Squared Errors (SSE) across different K values \citep{thorndikeWhoBelongsFamily1953}. The quality of the resulting clusters was validated using the Silhouette Coefficient \citep{rousseeuwSilhouettesGraphicalAid1987}. This process transformed the continuous MBT value of each measurement into a discrete cluster ID.

\subsubsection{Association Rule Mining for Weight Generation}\label{sec:apriori}
With the discretized data, we employed the Apriori algorithm to discover frequent co-occurrence patterns (itemsets) of MBT states across the 10 channels \citep{agrawalMiningAssociationRules1993a,hanMiningFrequentPatterns2000a}. This mining process was conducted independently for the earthquake-related and non-seismic datasets within each of the five environmental zones. To identify the most significant precursor signatures and quantitatively separate them from background variability, we used a **support-difference criterion**. This criterion, which calculates the difference in support (i.e., frequency of occurrence) for each itemset between the earthquake and non-seismic datasets, serves as a robust proxy for the signal-to-noise ratio. Itemsets with high positive support differences, indicating a statistically significant higher prevalence before earthquakes, were retained.

These results were then used to generate a quantitative weight for each frequency channel. For each channel, we calculated its average normalized support-difference across all significant itemsets in which it appeared. This data-driven weight reflects the channel's relative importance for precursor detection within a specific environment and forms the core of the domain knowledge integrated into our deep learning model.

\subsection{The WE-FTT Model for Anomaly Detection}\label{sec:weftt_model}
We selected a Transformer-based architecture due to its inherent ability to model complex, long-range dependencies. The self-attention mechanism is ideally suited to capture the synergistic interplay between different frequency and polarization channels, a characteristic identified as critical by our association rule mining.

\subsubsection{Weight-Enhanced Feature Processing}\label{sec:weight_processing}
Based on the mining results, we designed a feature processing approach to explicitly integrate the derived domain knowledge. Each of the 10 MBT feature columns was augmented with a corresponding weight column, whose value was pre-computed from the mining analysis (initial value of 1 plus the calculated support-difference weight). This creates a dual-input architecture where the model receives both the MBT features and their importance weights simultaneously. This explicit injection of channel importance guides the model to prioritize the most reliable features, reducing its susceptibility to noise from less informative channels. This process is illustrated in the parallel pathways of Figure~\ref{fig:fig12}, where the left pathway processes MBT features and the right pathway processes their corresponding weights before element-wise multiplication.

\subsubsection{Adaptive FT-Transformer Architecture}\label{sec:adaptive_architecture}
The core of our model is an adapted version of the FT-Transformer, which is designed for tabular data \citep{gorishniyRevisitingDeepLearning2021,vaswaniAttentionAllYou2017}. We implemented several critical modifications to optimize it for this task. Instead of a separate CLS token, our model learns a representation by combining feature embeddings with their corresponding projected weight embeddings through element-wise multiplication. Furthermore, we found that placing LayerNorm *after* positional encoding was optimal for preserving the structural characteristics of MBT data. The model's forward propagation begins with parallel projection pathways for the 10 MBT features and their 10 corresponding weights, transforming them into high-dimensional embeddings. These are then combined via element-wise multiplication, effectively creating a weighted feature space. This combined representation is then processed through a series of attention blocks, with the final representation passed to a prediction head for classification. The WE-FTT model contains approximately 12.3 million trainable parameters. Training on a single NVIDIA A100 GPU requires approximately 6 hours for 100 epochs, while inference processes 10,000 samples in under 2 seconds. The complete architecture is shown in Figure~\ref{fig:fig12}.

\subsubsection{Enhanced Multi-Head Attention Mechanism}\label{sec:attention_mechanism}
Within each attention block, a linear projection layer simultaneously generates the query (Q), key (K), and value (V) matrices from the input tensors. These are then divided among multiple attention heads. The attention scores are calculated using the standard scaled dot-product attention formula. Crucially, the weights' influence is already integrated into the input embeddings before the attention calculation. This indirect approach allows for more stable gradient flow while ensuring the attention mechanism operates on a representation where feature importance is already encoded. By focusing on these pre-weighted embeddings, the model can more effectively learn the complex inter-channel relationships crucial for distinguishing seismic precursors from environmental noise. The outputs from all heads are then concatenated and passed through a final projection layer to produce the output for the block.

\subsection{Model Training, Evaluation, and Uncertainty Quantification}\label{sec:model_evaluation}

\subsubsection{Optimization and Loss Function}\label{sec:optimization}
To ensure optimal and robust model training, we implemented a comprehensive optimization strategy. To address the class imbalance inherent in seismic data, we developed a Dynamic Focal Loss function. This function uses a gamma parameter (tunable between 0.5--5.0) to focus training on hard-to-classify samples and incorporates a momentum-based mechanism (momentum: 0.5--0.99) to adaptively adjust class weights based on recent performance. For hyperparameter optimization, we employed Optuna \citep{akibaOptunaNextgenerationHyperparameter2019} to conduct a systematic search over 30 trials with an early stopping criterion. Key optimized parameters included the number of attention heads (4--32), embedding dimension (64--512), number of attention blocks (2--8), learning rate (1e-5 to 1e-2), and weight decay (1e-6 to 1e-3). Training utilized a distributed multi-GPU setup, an 80/20 train-test split, and gradient clipping (0.5--5.0) to ensure stability.

\subsubsection{Performance Metrics and Comparative Analysis}\label{sec:metrics_comparison}

The model's effectiveness was evaluated using a suite of complementary metrics. The Matthews Correlation Coefficient (MCC) \citep{matthewsComparisonPredictedObserved1975} was chosen as the primary optimization metric due to its robustness in handling imbalanced multi-class problems. We also monitored accuracy, weighted precision-recall, F1 score, and Cohen's Kappa \citep{cohenCoefficientAgreementNominal1960} to provide a comprehensive assessment of performance.



To rigorously benchmark our WE-FTT model, we conducted a comparative analysis against several established machine learning models: RandomForest \citep{breiman2001random}, LightGBM, XGBoost \citep{chen2016xgboost}, TabNet \citep{arikTabNetAttentiveInterpretable2021}, CatBoost \citep{prokhorenkovaCatBoostUnbiasedBoosting2018}, and a Multi-Layer Perceptron (MLP). All models were trained and evaluated on identical dataset splits using the same input features (10 MBT channels and their 10 corresponding weight vectors). This extensive comparison framework allows for a robust assessment of the performance gains achieved by our specialized, knowledge-guided architecture. It is important to note that Principal Component Analysis (PCA), referenced later in Figure~\ref{fig:fig6}c, was used for visualization of the metric space only and was not treated as a predictive baseline.

\subsubsection{Uncertainty Quantification}\label{sec:uncertainty}
To assess the robustness of our methodological framework, we conducted several uncertainty analyses. Bootstrap resampling (1,000 iterations) was used to generate 95\% confidence intervals for the support values in the association rule mining, the results of which are visualized in our performance analyses. We performed sensitivity analysis on the K-means clustering by varying K from 3 to 20 and examining the stability of the resulting precursor signatures. Additionally, a stratified 5-fold cross-validation scheme was employed during model evaluation to ensure consistent performance across different data splits while maintaining the proportion of samples from each environmental zone.

\section{Results}
\subsection{Resolving the Signal Ambiguity Through Environmental Context}
Our analysis demonstrates that the signal ambiguity that has plagued satellite earthquake monitoring for decades can be effectively resolved through environment-specific analysis. By segmenting our global dataset into five distinct environmental zones and analyzing precursor patterns independently within each zone, we discovered that what appears as intractable noise in global analyses resolves into clear, consistent signatures when properly contextualized.

Figure~\ref{fig:fig2} demonstrated this breakthrough---each environmental zone exhibits unique, highly reliable frequency-polarization combinations for detecting seismic precursors. These environment-specific signatures achieve support values approaching or reaching 1.0, indicating near-perfect reliability within their respective environments. This finding fundamentally challenges the prevailing assumption that earthquake precursors should manifest uniformly across different surface types.

\subsection{Analysis of Multi-frequency Microwave for Zone-Dependent Pre-seismic Anomaly Detection}
An initial exploratory analysis across the full earthquake catalog (M$\ge$4.8) revealed that consistent, statistically significant MBT anomalies were predominantly associated with major earthquakes (M$\ge$7.0). For smaller magnitudes, potential signals were indistinguishable from background noise at a global scale, confirming that the signal ambiguity problem is particularly acute for lower-magnitude events. We therefore focused our main analysis on M$\ge$7.0 earthquakes, where signals, though still ambiguous in conventional global analyses, possess sufficient strength to become distinguishable through our environment-specific framework. The integrated methodological framework presented in Figure~\ref{fig:fig1} enabled comprehensive analysis of AMSR-2 microwave brightness temperature data across five distinct environmental zones has revealed characteristic pre-seismic anomaly patterns through specific frequency-polarization combinations, demonstrating significant detection capabilities for seismic precursor signals. Through comprehensive frequent itemset mining analysis, we have identified environment-specific anomaly signatures that achieve perfect or near-perfect support values, indicating robust seismic precursor detection capabilities across diverse landscape conditions (Figure~\ref{fig:fig2}).

In marine environments (Zone A), distinct pre-seismic anomalies were most effectively captured by the combination of 89~GHz H-polarization (anomaly range $\sim$219--253~K) and 36.5~GHz V-polarization (123--248~K), achieving perfect detection reliability (support = 1.0000). A secondary anomaly signature comprising 23.8~GHz H-polarization (177--235~K) and 36.5~GHz H-polarization (109--209~K) also showed near-perfect reliability (support = 0.9923). The overlapping temperature range ($\sim$177--209~K) between these channels defines a critical detection window for marine earthquake precursors.

Humid forest zones (Zone B) exhibited unique anomaly signatures, with optimal detection achieved through 89~GHz V-polarization (98--295~K) combined with 36.5~GHz V-polarization (140--297~K), yielding perfect support. Detection reliability showed a systematic slight decrease when additional frequency channels were included: for example, including the 10.65~GHz H-polarization channel (169--385~K) maintained a support of 0.5154, while adding 6.9~GHz H-polarization (164--388~K) gave 0.5151. This minimal support difference ($\Delta$support $\approx$ 0.0003) quantifies the effect of dense vegetation on precursor signal detectability.

Dry forest environments (Zone C) revealed seismic precursor signatures through 36.5~GHz H-polarization (121--207~K) combined with 10.65~GHz V-polarization (160--220~K), achieving near-perfect detection reliability. The stability of the 36.5~GHz H-polarization anomaly range (standard deviation <0.1~K) indicates a consistent precursor signature. Secondary combinations including 6.9~GHz V-polarization (152--216~K) also achieved very high reliability (support = 0.9976), with anomalies concentrated around 159--207~K.

Wetland zones (Zone D) demonstrated complex anomaly patterns. Optimal detection involved various combinations incorporating 6.9~GHz H-polarization (73--166~K). Anomaly signatures showed a stratified pattern: low-frequency H-polarization exhibited broad detection ranges ($\Delta T \approx 93$ K), mid-frequency V-polarization had intermediate detection windows (e.g., $\Delta T \approx 164$ K for 23.8~GHz), and high-frequency channels provided complementary detection characteristics (e.g., $\Delta T \approx 214$ K for 36.5~GHz V-polarization). This stratification suggests multi-layered precursor processes in wetlands.

In arid zones (Zone E), seismic precursor detection was optimized through the combination of 23.8GHz H-polarization (189.04-240.88K) and 36.5GHz V-polarization (103.18-246.57K), achieving perfect reliability while individual channels showed significantly lower detection capabilities (0.0693 for 23.8GHz V). Low-frequency channels exhibited broad anomaly detection ranges ($\Delta T \approx 154$K for 6.9GHz H) despite limited reliability ($\le$0.005), indicating complex subsurface precursor mechanisms.

\subsection{Terrain-Specific Microwave Sensitivity and Optimal Channel Combinations}
To evaluate AMSR-2's multi-frequency microwave sensitivity to seismic anomalies, we conducted a comprehensive analysis of mean support values across diverse terrain types. The study categorized surface types into five zones based on vegetation coverage and soil moisture content: marine areas (Zone A), high vegetation-high soil moisture areas (Zone B), high vegetation-low soil moisture areas (Zone C), low vegetation-high soil moisture areas (Zone D), and low vegetation-low soil moisture areas (Zone E). We analyzed horizontal and vertical polarization channels at 6.9~GHz, 10.65~GHz, 23.8~GHz, 36.5~GHz, and 89.0~GHz frequencies, calculating mean support values across all frequent itemsets to assess seismic anomaly detection capabilities.

Analysis of mean support values revealed distinctive patterns across frequency bands and polarizations (see Table~\ref{tab:table1} and Figure~\ref{fig:fig3}). At 6.9~GHz, horizontal polarization demonstrated varying effectiveness across terrain types, with highest support values in Zone D (0.9348, MBT: 166.36-655.35K), followed by Zone C (0.5716, MBT: 73.61-145.63K). Vertical polarization at this frequency showed similar patterns but with generally higher MBT ranges, particularly in Zone D (0.9219, MBT: 214.9-655.35K). Zone B showed moderate support values (0.2688, MBT: 163.66-388.18K), while Zone A exhibited minimal sensitivity (0.0074, MBT: 41.88-168.06K).

The 10.65~GHz band exhibited enhanced detection capability, especially in wetland areas (Zone D), where horizontal polarization achieved a support value of 0.9306 (MBT: 171.55-655.35K) and vertical polarization reached 0.9229 (MBT: 219.56-655.35K). Notably, dry forest regions (Zone C) showed distinct polarization preferences, with vertical polarization achieving 0.6220 (MBT: 159.63-220.37K). Zones A and B maintained relatively low support values, with Zone B showing slightly better performance (0.2689, MBT: 168.79-385.17K).

Higher frequency bands demonstrated more pronounced environment-specific responses. The 23.8~GHz horizontal polarization showed exceptional performance in Zone D (0.9044, MBT: 201.73-655.35K), while maintaining moderate effectiveness in Zone E (0.5075, MBT: 189.04-240.88K). Zone B showed improved response at this frequency, reaching 0.3016 (MBT: 104.42-182.47K). The 36.5~GHz band achieved the highest overall support values, particularly in Zone D where horizontal polarization reached 0.9713 (MBT: 209.48-655.35K) and vertical polarization achieved 0.8738 (MBT: 139.48-301.56K). The 89.0~GHz band showed unique capabilities in arid regions (Zone E), with horizontal polarization achieving 0.7380 (MBT: 64.46-215.3K), while Zone A reached its maximum sensitivity at this frequency (0.5078, MBT: 219.34-252.97K).

Multiple frequency combinations demonstrated enhanced detection capabilities in specific terrains. In Zone D, the combination of 36.5~GHz and 23.8~GHz horizontal polarization channels provided optimal detection with support values consistently above 0.90. Zone C showed best response to combined 36.5~GHz horizontal polarization (0.7028, MBT: 121.4-207.09K) and 10.65~GHz vertical polarization (0.6219, MBT: 152.28-215.62K) observations. Zone E exhibited optimal sensitivity with 89.0~GHz horizontal polarization (0.7380, MBT: 64.46-215.3K) paired with 23.8~GHz vertical polarization (0.5224, MBT: 232.9-261.39K). Zones A and B showed limited improvement with frequency combinations, suggesting fundamental physical limitations in these environments.

\subsection{Model Performance Evaluation}
The proposed Weight-Enhanced Feature-Tailored Transformer (WE-FTT) architecture was systematically evaluated through comprehensive classification analysis and comparative performance assessment against established machine learning models. This section presents the detailed results of these evaluations, demonstrating the effectiveness of the frequency itemset-based weight optimization approach.

\subsubsection{Classification Performance of the WE-FTT Model}
The classification capabilities of the WE-FTT model were evaluated using a confusion matrix analysis across all frequency-polarization combinations, as shown in Figure~\ref{fig:fig4}. The model demonstrates remarkable classification accuracy, achieving an average of 84.2\% correct classifications across all microwave channels. This high accuracy is particularly significant given the complex nature of the classification task, which involves distinguishing between multiple frequency-polarization combinations in the presence of environmental variability and potential seismic anomalies.

The confusion matrix reveals several important patterns in the model's classification behavior. First, the diagonal elements (navy blue) show consistently high accuracy across all channels, with particularly strong performance in the 36.5~GHz (84.2\%) and 89.0~GHz (84.2%) bands. Second, the light blue cells, representing H-V polarization confusion within the same frequency, show minimal misclassification (average 3.1%), indicating the model's strong capability to distinguish between polarization states. This is particularly important for seismic anomaly detection, as polarization sensitivity provides critical information about surface structural changes associated with seismic activity.

The coral-colored cells represent cross-frequency confusion, with color intensity corresponding to the percentage of misclassifications. Notably, the confusion predominantly occurs between adjacent frequency bands (e.g., between 6.9~GHz and 10.65~GHz), which is physically reasonable given the spectral proximity of these channels. The model shows minimal confusion between spectrally distant channels (e.g., between 6.9~GHz and 89.0~GHz), with an average cross-band confusion of only 1.9%. This selective confusion pattern demonstrates the model's ability to capture the physical relationships between frequency channels while maintaining discriminative power between spectrally distinct measurements.

\subsubsection{Comparative Analysis of Model Performance}
To rigorously evaluate the effectiveness of the proposed WE-FTT approach, we conducted a comprehensive performance comparison against five established machine learning models: RandomForest, LightGBM, TabNet, CatBoost, and XGBoost. Each model was trained on identical dataset splits and evaluated using six complementary performance metrics: Matthews Correlation Coefficient (MCC), F1 Score, Accuracy, Precision, Cohen's Kappa, and Recall. Figure~\ref{fig:fig5} presents the detailed comparison across these metrics, revealing consistent performance advantages for the WE-FTT architecture.

As shown in Figure~\ref{fig:fig5}, the proposed WE-FTT model achieved superior performance across all evaluation metrics. The WE-FTT achieved an MCC of 0.84, F1 score of 0.82, accuracy of 0.84, precision of 0.80, Cohen's Kappa of 0.82, and recall of 0.84. These results represent substantial improvements over the next best model (RandomForest), which achieved an MCC of 0.74, F1 score of 0.70, accuracy of 0.72, precision of 0.71, Cohen's Kappa of 0.68, and recall of 0.72. This performance gap is especially significant considering that MCC is particularly sensitive to imbalanced classification scenarios characteristic of seismic anomaly detection tasks. The consistent performance advantages across diverse metrics demonstrate the WE-FTT's robust generalization capabilities across different evaluation criteria. The significantly higher MCC score, in particular, highlights the WE-FTT's enhanced capability in handling the imbalanced nature and inherent signal ambiguity of the seismic precursor detection task, where baseline models likely struggle more to differentiate true positives from complex background noise.

Statistical significance testing, as illustrated in Figure~\ref{fig:fig6}a, confirms the reliability of these performance differences. Bootstrap resampling with 1,000 iterations generated 95% confidence intervals for the aggregated performance scores, revealing statistically significant differences between the WE-FTT (0.8271) and RandomForest (0.7119, p<0.01), as well as between subsequent model pairs in the performance ranking. This statistical validation reinforces the conclusion that the WE-FTT's performance advantages are systematic rather than artifacts of sampling variation.

The multi-dimensional performance profile visualization in Figure~\ref{fig:fig6}b further illustrates the WE-FTT's balanced excellence across all metrics. While some baseline models showed uneven performance profiles with strengths in certain metrics and weaknesses in others, the WE-FTT maintained consistently high values across all dimensions, creating a more uniformly expanded performance polygon. This balanced profile indicates that the model achieves its high overall performance without sacrificing specific aspects of classification quality.

Principal Component Analysis (PCA) of the performance metrics, shown in Figure~\ref{fig:fig6}c, provides additional insights into the relationships between models and evaluation criteria. The first two principal components capture 93.5% of the total variance, with PC1 (78.2%) primarily representing overall performance and PC2 (15.3%) capturing the precision-recall trade-off dimension. The WE-FTT's position in the upper-right quadrant of this reduced feature space demonstrates its dominance in both dimensions, while the clustering of lower-performing models in the left region highlights the significant performance gap between the proposed approach and traditional methods.

To further examine the detailed classification patterns across different models, Figure~\ref{fig:fig7} presents the confusion matrices for all six evaluated models. The WE-FTT model (Figure~\ref{fig:fig7}a) demonstrates substantially higher diagonal values (correct classifications) compared to all other models, with an overall accuracy of 84.2%. In contrast, the next best performer, RandomForest (Figure~\ref{fig:fig7}b), achieves only 71.9% accuracy, with notably increased cross-frequency confusion particularly between adjacent frequency bands. The performance differential becomes even more pronounced when comparing with the lower-performing models: LightGBM (67.6%), TabNet (63.0%), CatBoost (58.1%), and XGBoost (55.4%).

A particularly noteworthy pattern emerges when examining the H-V polarization confusion (light blue cells) across different models. The WE-FTT demonstrates remarkably low in-band polarization confusion, while traditional models show significantly higher confusion rates in this category. This enhanced polarization discrimination capability can be directly attributed to the weight optimization strategy, which effectively leverages the distinctive information content of different polarization channels based on their relevance to seismic anomaly detection.

The confusion matrices also reveal that lower-performing models exhibit increased misclassification between spectrally distant channels. This pattern suggests that traditional models struggle to capture the complex physical relationships between frequency channels, while the WE-FTT's attention mechanism effectively models these interdependencies through its weight-enhanced feature interactions.

\subsubsection{Relationship to Frequency Itemset Mining Results}
The superior performance of the WE-FTT model can be directly linked to the frequency itemset mining results presented earlier. The model's pre-computed weight approach effectively leverages the identified optimal frequency-polarization combinations for each environmental zone. These weights, derived from the support values identified in the itemset mining analysis, are processed through a dedicated projection pathway before being combined with feature embeddings through element-wise multiplication. For instance, in Zone D (Wetland), the model assigns higher weights to the 36.5~GHz and 23.8~GHz horizontal polarization channels, which were identified as having exceptionally high support values (0.9713 and 0.9044 respectively) in the itemset mining analysis. Similarly, for Zone C (Dry Forest), the model emphasizes the 36.5~GHz horizontal polarization (support value: 0.7028) and 10.65~GHz vertical polarization (support value: 0.6220) channels.

This superior ability to resolve ambiguous signals stems directly from integrating the mined knowledge. By assigning pre-computed weights based on rules identifying consistent pre-seismic signatures (e.g., high support for 36.5~GHz H and 10.65~GHz V in Zone C), the WE-FTT effectively learns the subtle, environment-specific decision boundary between seismic and non-seismic states, enabling more reliable classification even when individual channel readings might otherwise be misleading. By focusing attention on the most informative frequency-polarization combinations for each environmental context, the WE-FTT achieves both improved classification accuracy and enhanced physical interpretability of the learned patterns.

The significant performance improvements observed in the comparative analysis thus validate the effectiveness of the frequency itemset-based weight optimization methodology. By systematically mining the relationships between frequency channels and their association with seismic activity across different environmental conditions, and then integrating this knowledge into the model's attention mechanism, the proposed approach successfully addresses the complex challenges of microwave-based seismic anomaly detection.


\subsection{Ablation Study}
To understand the mechanistic basis of our model's superior performance, we systematically evaluated the contribution of each architectural component through comprehensive ablation experiments. This analysis not only validates our design choices but also reveals fundamental insights into how different frequency channels contribute to seismic precursor detection.

\subsubsection{Architectural Component Hierarchy}
Our ablation experiments across 18 variants revealed a clear hierarchy of component importance (Figure~\ref{fig:fig8}). The weight-enhanced dual projection pathway emerged as the cornerstone of our architecture---its removal caused catastrophic performance collapse (MCC: 0.84$\rightarrow$0.42), representing nearly 50\% degradation. This dramatic impact confirms that pre-computed frequency importance weights cannot be effectively learned through end-to-end training alone. Feature projection proved similarly critical (41.4\% impact), while attention mechanisms showed more nuanced effects: residual connections (35\% impact) outweighed multi-head attention benefits (28\% impact), suggesting that stable gradient flow trumps attention diversity for geophysical signal processing. Surprisingly, position encoding---fundamental in language models---contributed minimally (<6\%), indicating that spatial relationships in gridded MBT data differ fundamentally from sequential dependencies.

\subsubsection{Performance Signatures Across Metrics}
Different architectural modifications produced distinct performance signatures across evaluation metrics (Figure~\ref{fig:fig9}). Weight projection removal caused asymmetric degradation: Cohen's Kappa plummeted 57.7\% while Precision dropped only 44.5\%, revealing that weight integration specifically enhances cross-class discrimination---crucial for distinguishing genuine precursors from environmental noise. Feature fusion variants maintained more balanced profiles, with alternative fusion achieving near-baseline MCC (0.816) while selectively compromising F1 Score (0.755). This metric-specific sensitivity provides a roadmap for future architectural refinements: components showing uniform degradation (weight projection) are irreplaceable, while those with selective impacts (fusion strategies) offer optimization opportunities.

\subsubsection{Component-Channel Coupling}
Integrated analysis uncovered a profound connection between architectural components and physical signal characteristics (Figure~\ref{fig:fig10}). Component impact quantification (Figure~\ref{fig:fig10}a) established that input projection mechanisms account for over 70\% of total model performance, far exceeding traditional attention mechanisms. Multi-metric analysis (Figure~\ref{fig:fig10}b) revealed specialized roles: input projections govern overall accuracy, attention mechanisms control error distributions, while fusion strategies balance class predictions. The effectiveness-complexity trade-off (Figure~\ref{fig:fig10}c) identified unexpected efficiency gains---learnable position embeddings slightly improved performance while reducing computational overhead, challenging conventional transformer design wisdom.
Most revealing was the class-specific sensitivity analysis (Figure~\ref{fig:fig10}d), which demonstrated remarkable alignment between architectural importance and physical significance. Frequency channels 0, 3, 4, and 9---precisely those identified as optimal for seismic detection in our mining analysis---showed 3-fold higher sensitivity to weight projection removal. This direct correspondence validates our central hypothesis: integrating domain knowledge through architectural design enables detection of subtle geophysical patterns invisible to conventional approaches.

\subsubsection{Implications for Geophysical Signal Processing} 
These findings establish three fundamental principles for AI-driven geophysical analysis: First, domain knowledge integration through specialized architectural pathways outperforms generic deep learning. Second, the hierarchy of component importance mirrors the physical hierarchy of signal generation---from lithospheric stress to surface emissions. Third, the tight coupling between mined frequency patterns and architectural sensitivity suggests that our model has learned physically meaningful representations rather than statistical artifacts.
The ablation study thus transcends technical validation, revealing how architectural innovations can encode geophysical understanding into neural networks. By demonstrating that pre-computed weights based on seismic-frequency associations are irreplaceable by learned parameters, we establish a new paradigm for knowledge-guided deep learning in Earth observation.

\subsection{Global validation}\label{sec:global_validation}
To address the critical need for robust, global-scale validation, we conducted two supplementary experiments as requested by the reviewers. These tests were designed to (1) quantify the model's false-positive rate (FPR) on non-seismic days across the globe and (2) evaluate its performance on recent, out-of-sample earthquake events.

First, we performed a global random-date test by applying our trained model to over 100 randomly selected days for each of the five environmental zones, ensuring no overlap with any known seismic event windows. The resulting per-pixel FPR is visualized in a world map (Fig.~\ref{fig:fig13}), with zonal statistics summarized in Table~\ref{tab:tableS1}. This analysis provides a clear baseline for the model's performance in the absence of seismic signals.

Second, we conducted an out-of-sample validation test on all major (M$\ge$7.0) earthquakes that occurred after our dataset's temporal cutoff (post-August 2023). Using the frozen, pre-trained WE-FTT model, we performed a forward-looking test. The performance on these new events, reported in Table~\ref{tab:tableS2} and illustrated in Figs.~\ref{fig:figS1}-\ref{fig:figS2}, demonstrates the model's temporal robustness and generalization capability.

Finally, a dedicated sensitivity analysis was performed to explicitly control for potential artifacts from tsunamis and extreme sea states in marine environments. By applying strict exclusion windows, we confirmed that significant pre-seismic detections persist, as detailed in Table~\ref{tab:tableS3} and Fig.~\ref{fig:figS3}. Together, these validations provide strong evidence for the reliability and robustness of our environment-specific framework.

\section{Discussion}\label{sec:discussion}

\subsection{Resolving Signal Ambiguity via an Environment-Specific Framework}\label{sec:discuss_ambiguity}

Our findings demonstrate a robust pathway to resolving the signal ambiguity problem that has challenged satellite-based earthquake monitoring for decades. The key insight, supported by the near-perfect support values for specific channel combinations in distinct environmental zones (Table~\ref{tab:table1}), is that reliable pre-seismic signatures in MBT data are not universal but are fundamentally environment-specific. By abandoning the search for a single, global precursor pattern and instead embracing environmental heterogeneity, our framework successfully transforms what appears as intractable noise in global analyses into clear, context-dependent indicators of potential seismic preparation. This paradigm shift is validated by the superior performance of the WE-FTT model (MCC of 0.84), which significantly outperforms conventional global methods by effectively learning these nuanced, zone-dependent signatures.

\subsection{Physical Interpretation and Comparison with Prior Work}\label{sec:discuss_physical}

\subsubsection{Environment-Specific Physical Mechanisms}
The observed frequency-dependent anomaly patterns provide compelling evidence for multiple, interconnected physical mechanisms. In wetland and arid zones (Zones D and E), the high support for low-frequency channels (e.g., 6.9 GHz H-pol) aligns with the P-hole activation hypothesis, where stress-induced positive charge carriers modify subsurface dielectric properties \citep{freundPreearthquakeSignalsUnderlying2011}. The deeper penetration of these frequencies is crucial for detecting such phenomena. Conversely, marine environments (Zone A) show distinct signatures in high-frequency channels (e.g., 89 GHz H-pol), consistent with thermal energy transfer or microscale surface roughness changes from submarine activity \citep{liuGeneralFeaturesMultiparameter2023}.

\subsubsection{Signal Discrimination from Background Noise and Contrast with Prior Studies}
Our framework represents a significant methodological advance over prior studies that often relied on single-channel analysis or fixed anomaly thresholds \citep{wuIdentifyingSeismicAnomalies2024}. Our results (Fig.~\ref{fig:fig2}) prove that certain channel combinations achieve near-perfect reliability *only when analyzed together*, indicating non-linear inter-channel synergies. Our support-difference criterion, detailed in \S\ref{sec:data_identification}, serves as a quantitative proxy for the signal-to-noise ratio, statistically separating precursor patterns from background variability. Furthermore, our work contrasts with ionospheric precursor studies, such as the analysis of Total Electron Content (TEC). For instance, \citet{Shah2020TotalEC} identified TEC anomalies with varying lead times and magnitudes, noting dependencies on earthquake depth and latitude. While TEC anomalies reflect ionospheric perturbations, our MBT anomalies are linked to lithosphere-atmosphere interactions. Our new seismological stratification analysis (\S\ref{sec:seismo_controls}) allows for a more direct comparison, revealing that, similar to TEC, MBT anomaly characteristics also show dependence on source parameters, reinforcing the concept of a multi-layered Lithosphere-Atmosphere-Ionosphere Coupling (LAIC) system.

\subsection{Methodological Contributions and Physics-Informed AI}\label{sec:discuss_implications}

This work introduces a novel paradigm for Earth observation by architecturally encoding domain knowledge into a deep learning model. The ablation study (Figs.~\ref{fig:fig8}-\ref{fig:fig10}) confirms that the weight-enhanced dual projection pathway, which integrates our mined frequency-importance rules, is the most critical component, with its removal causing a catastrophic 50\% performance drop. This demonstrates that pre-computed weights based on physical associations cannot be effectively replaced by end-to-end learning alone, establishing a new blueprint for "physics-informed AI" in the geosciences.

\subsection{Limitations and Future Directions}\label{sec:discuss_limitations}

Despite these advances, several limitations remain. Our analysis focuses on M$\ge$7.0 earthquakes, and the 0.25$^{\circ}$ spatial resolution may be too coarse for smaller events or highly localized phenomena. Methodologically, the K-means clustering and Dobrovolsky radius are simplifications. Potential confounding from extreme weather, volcanic activity, or anthropogenic sources remains a challenge.

Future work should focus on three key areas. First, developing region-specific models to test applicability to moderate-magnitude earthquakes. Second, integrating multi-parameter observations (e.g., GNSS strain, gas concentrations, TEC) for more robust, multi-faceted validation of anomalies \citep{hayakawaIntegratedAnalysisMultiParameter2024}. Third, real-time implementation of this framework, leveraging the identified optimal channels and adaptive, environment-specific thresholds, is a critical next step toward operational deployment.


\section{Conclusion}

This study successfully addresses the long-standing signal ambiguity problem in satellite-based earthquake monitoring through the development and validation of an environment-specific, knowledge-guided analytical framework. Our analysis of 346.56 million AMSR-2 observations for 154 major (M$\ge$7.0) earthquakes demonstrates that reliable pre-seismic MBT signatures are not universal but manifest distinctly across five environmental zones, with specific frequency-polarization combinations achieving near-perfect detection reliability. This finding explains the limited success of previous global, uniform approaches.

Our key methodological contribution is the Weight-Enhanced Feature-Tailored Transformer (WE-FTT), a novel architecture that architecturally integrates domain knowledge of these environment-specific signatures. This "physics-informed AI" approach achieved a Matthews Correlation Coefficient of 0.84, substantially outperforming conventional methods and confirming that accounting for environmental heterogeneity is a fundamental requirement for robust precursor detection. This research therefore provides both a practical solution to the signal ambiguity problem and a generalizable framework for other remote sensing applications where environmental context fundamentally modulates signal characteristics.

Looking forward, the integration of this framework with complementary observations, such as GNSS-derived strain and ionospheric perturbations, represents the next critical step toward building robust, multi-parameter monitoring systems. Real-time implementation using the identified optimal channels for each environmental zone will be essential for advancing toward operational deployment.
% ===== Figures =====
\clearpage
\begin{figure*}[p]
\centering
\includegraphics[width=1.0\linewidth]{image1.png}
\caption{Resolving signal ambiguity through environment-specific analysis: methodological framework. The workflow illustrates the integration of surface type classification, earthquake-related data sampling, clustering analysis, association rule mining, and the WE-FTT model.}\label{fig:fig1}
\end{figure*}
\clearpage

\begin{figure*}[!t]
\centering
\includegraphics[width=1.1\linewidth]{image11.png}
\caption{Global Distribution of Major Earthquakes (M$\ge$7.0) Across Different Surface Type Zones}\label{fig:fig11}
\end{figure*}
\clearpage

\begin{figure*}[p]
\centering
\includegraphics[width=1.1\linewidth]{image12.png}
\caption{Architectural overview of the WE-FTT.}\label{fig:fig12}
\end{figure*}
\clearpage

\begin{figure*}[p]
\centering
\includegraphics[width=1.0\linewidth]{image2.png}
\caption{Environment-specific precursor signatures achieving near-perfect reliability.}\label{fig:fig2}
\end{figure*}
\clearpage

\begin{figure*}[p]
\centering
\includegraphics[width=1.1\linewidth]{image3.png}
\caption{Polarization dependence of microwave response characteristics across terrain zones.}\label{fig:fig3}
\end{figure*}
\clearpage

\begin{figure*}[p]
\centering
\includegraphics[width=1.1\linewidth]{image4.png}
\caption{Brightness Temperature Polarization Classification Matrix across five AMSR-2 frequency bands.}\label{fig:fig4}
\end{figure*}
\clearpage

\begin{figure*}[p]
\centering
\includegraphics[width=1.1\linewidth]{image5.png}
\caption{Model performance comparison across six evaluation metrics.}\label{fig:fig5}
\end{figure*}
\clearpage

\begin{figure*}[p]
\centering
\includegraphics[width=1.1\linewidth]{image6.png}
\caption{Comprehensive model performance analysis and PCA visualization. (a) Statistical significance of performance differences with 95% confidence intervals from bootstrap resampling. (b) Multi-dimensional performance profiles across six key metrics. (c) Principal Component Analysis (PCA) of the performance metrics, used here for visualization of the metric space only; it is not a predictive baseline.}\label{fig:fig6}
\end{figure*}
\clearpage

\begin{figure*}[p]
\centering
\includegraphics[width=1.0\linewidth]{image7.png}
\caption{Confusion matrices showing classification performance across different models.}\label{fig:fig7}
\end{figure*}
\clearpage

\begin{figure*}[p]
\centering
\includegraphics[width=1.1\linewidth]{image8.png}
\caption{Ablation study results across key architectural components.}\label{fig:fig8}
\end{figure*}
\clearpage

\begin{figure*}[p]
\centering
\includegraphics[width=1.1\linewidth]{image9.png}
\caption{Radar visualization of model performance metrics across ablation variants.}\label{fig:fig9}
\end{figure*}
\clearpage

\begin{figure*}[p]
\centering
\includegraphics[width=1.1\linewidth]{image10.png}
\caption{Comprehensive ablation analysis of the WE-FTT model.}\label{fig:fig10}
\end{figure*}
\clearpage





% ===== Tables =====
\clearpage
\begin{table*}[p]
\small
\centering
\caption{Segmented MBT ranges and mean support values across terrain zones.}\label{tab:table1}
\resizebox{\textwidth}{!}{%
 \renewcommand{\arraystretch}{1.1}
\begin{tabular}{@{}llccccc@{}}
\toprule
      &   & \multicolumn{5}{c}{Mean Support Value (MBT segmentation range)} \\
 \cmidrule(lr){3-7}
\multicolumn{2}{l}{Freq. (GHz)}  & Zone A & Zone B & Zone C & Zone D & Zone E \\
\midrule
\multirow{4}{*}{6.9}
    & \multirow{2}{*}{H}
        & 0.0074\,(41.8--168.1)  & 0.2688\,(163.7--388.2) & 0.5716\,(73.6--145.6)  & 0.9348\,(166.4--655.4) & 0.4038\,(168.9--655.4) \\
    &   & ---                   & 0.2201\,(73.4--163.7)  & 0.4408\,(226.8--382.8) & 0.8015\,(73.3--166.4)  & 0.0042\,(14.3--168.9) \\[6pt]
    & \multirow{2}{*}{V}
        & 0.0076\,(57.1--219.2)  & \textbf{0.3095\,(252.2--388.1)} & 0.6219\,(152.3--215.6) & 0.9219\,(214.9--655.4) & 0.4038\,(218.6--655.4) \\
    &   & ---                   & 0.2193\,(138.1--205.0) & 0.4743\,(215.6--384.4) & 0.7759\,(139.7--214.9) & 0.0043\,(47.4--218.6) \\
\midrule
\multirow{4}{*}{10.65}
    & \multirow{2}{*}{H}
        & 0.0075\,(68.9--173.7)  & 0.2689\,(168.8--385.2) & 0.5578\,(79.0--130.8)  & 0.9306\,(171.6--655.4) & 0.4038\,(173.4--655.4) \\
    &   & ---                   & 0.2202\,(78.9--168.8)  & ---                   & 0.8014\,(78.9--171.5)  & 0.0041\,(79.1--173.4) \\[6pt]
    & \multirow{2}{*}{V}
        & 0.0075\,(121.3--223.9) & 0.2198\,(146.2--212.0) & 0.6220\,(159.6--220.4) & 0.9229\,(219.6--655.4) & 0.4038\,(222.8--655.4) \\
    &   & ---                   & ---                   & 0.4748\,(220.4--385.4) & 0.7690\,(148.0--219.6) & 0.0042\,(146.6--222.8) \\
\midrule
\multirow{5}{*}{23.8}
    & \multirow{3}{*}{H}
        & \textbf{0.5124\,(177.2--234.6)} & 0.3016\,(104.4--182.5) & 0.4676\,(233.8--292.5) & 0.9044\,(201.7--655.4) & 0.5075\,(189.0--240.9) \\
    &   & 0.0132\,(103.3--177.2) & 0.2506\,(237.0--292.1) & 0.2179\,(180.3--233.8) & ---                   & 0.3366\,(104.8--189.0) \\
    &   & ---                   & 0.1465\,(182.5--237.0) & ---                   & ---                   & --- \\[6pt]
    & \multirow{2}{*}{V}
        & 0.0183\,(241.0--655.3) & 0.3029\,(157.2--256.0) & 0.5043\,(257.4--302.6) & 0.8738\,(201.7--655.4) & 0.5224\,(232.9--261.4) \\
    &   & 0.0069\,(162.9--241.0) & 0.1549\,(256.0--298.3) & 0.0497\,(227.9--257.4) & ---                   & 0.0426\,(168.1--232.9) \\
\midrule
\multirow{4}{*}{36.5}
    & \multirow{2}{*}{H}
        & 0.0227\,(109.3--209.3) & 0.1745\,(121.7--185.1) & \textbf{0.7028\,(121.4--207.1)} & \textbf{0.9713\,(209.5--655.4)} & 0.4052\,(210.9--655.4) \\
    &   & ---                   & ---                   & 0.5494\,(207.1--290.9) & 0.7824\,(111.5--209.5) & 0.0004\,(102.5--210.9) \\[6pt]
    & \multirow{2}{*}{V}
        & 0.0228\,(122.7--247.6) & 0.3054\,(139.6--297.2) & 0.5412\,(148.4--244.3) & 0.8738\,(139.5--301.6) & 0.4079\,(103.2--246.6) \\
    &   & ---                   & ---                   & 0.4229\,(244.3--300.6) & ---                   & 0.0305\,(246.6--655.4) \\
\midrule
\multirow{5}{*}{89.0}
    & \multirow{3}{*}{H}
        & 0.5078\,(219.3--253.0) & 0.3053\,(94.8--212.0)  & 0.5567\,(217.3--250.0) & 0.8738\,(80.6--295.9)  & \textbf{0.7380\,(64.5--215.3)}  \\
    &   & 0.0261\,(253.0--644.4) & 0.1638\,(246.1--292.9) & 0.0340\,(250.0--293.4) & ---                   & 0.0159\,(215.3--250.0) \\
    &   & 0.0023\,(72.6--219.3)  & 0.1335\,(212.1--246.1) & ---                   & ---                   & ---\\[6pt]
    & \multirow{2}{*}{V}
        & 0.0179\,(270.7--655.4) & 0.3054\,(98.0--295.4)  & 0.5741\,(260.8--298.5) & 0.8738\,(82.8--298.8)  & 0.2657\,(65.2--231.2)  \\
    &   & 0.0018\,(254.1--570.7) & ---                   & 0.0001\,(229.1--260.8) & ---                   & 0.0647\,(231.3--263.3) \\
\bottomrule
\end{tabular}}
\end{table*}

\begin{table*}[p]
\small
\centering
\caption{Classification criteria for different surface types}\label{tab:table2}
\resizebox{\textwidth}{!}{%
 \renewcommand{\arraystretch}{1.1}
\begin{tabular}{@{} c c c c l @{}}
\toprule
\textbf{Zone} & \textbf{Surface Type} & \textbf{Vegetation Coverage} & \textbf{Soil Moisture} & \textbf{Characteristics} \\
\midrule
A & Marine Zone & N/A & N/A & Oceanic regions \\
B & Humid Forest & $>0.5$ & $>6$ & Tropical rainforests, dense vegetation with moist soil \\
C & Dry Forest & $>0.5$ & $\le 6$ & Subtropical savannas, significant vegetation with dry soil \\
D & Wetland & $\le 0.5$ & $>6$ & Wetlands and marshes \\
E & Arid Land & $\le 0.5$ & $\le 6$ & Deserts and arid areas \\
\bottomrule
\end{tabular}}
\end{table*}

% ===== Back matter and references =====
\clearpage
\nolinenumbers\bibliographystyle{elsarticle-harv}
\bibliography{references}

\section*{Acknowledgements}
We thank the data providers for making their datasets publicly available. This work was supported by the Natural Science Foundation of China (NSFC) under grant number 42274108 and the Central Public-interest Scientific Institution Basal Research Fund (No. CEAIEF2025030105).

\section*{Author contributions statement}
P.X., C.L., and X.S. conceived the study. P.X. and C.L. performed the analysis and modelling. H.Z., R.B., A.D.S and X.S. supervised the work, contributed to interpretation and discussion. P.X. wrote the initial manuscript. All authors discussed the results and reviewed the manuscript.

\section*{Additional information}
\textbf{Data Availability} All data supporting the findings of this study are publicly available. The AMSR-2 brightness temperature datasets (2013--2023) analyzed in this study can be accessed through the JAXA GCOM-W1 data archive (\url{https://gcom-w1.jaxa.jp/}). Daily ERA5-Land reanalysis data for vegetation cover are available from the ECMWF data portal (\url{https://www.ecmwf.int/}), and GLDAS soil moisture data were obtained from NASA's HDISC (\url{https://disc.gsfc.nasa.gov/}). The global earthquake catalog was downloaded from the USGS Earthquake Hazards Program (\url{https://www.usgs.gov/programs/earthquake-hazards/earthquakes}). Processed datasets (e.g., discretized time series and derived association rules) are available from the corresponding author upon reasonable request.

\textbf{Code Availability} The PyTorch implementation of WE-FTT, along with all code for data preprocessing, model training, and analysis, is publicly available on GitHub at \url{https://github.com/PANXIONG-CN/WE-FTT} under the MIT License.

\textbf{Competing interests} The authors declare no competing interests.

\end{document}
